\section{Results}\label{sec:results}

\subsection{Aggregate randomized-task performance}

Table \ref{tab:primary-results} and Figure \ref{fig:primary-results} report the
primary randomized task law.  The source-scored design contains a truly
feasible policy in 91.7\% of task--resolution cells and yields a certified
true-feasible deployment in 47.3\%.  In the supplementary comparisons,
standardized DCT maximin reaches 81.0\% and 31.5\%, and the source-greedy
portfolio reaches 87.1\% and 37.1\%.  Penalized losses are 0.163, 0.274, and
0.178, respectively.  Source scoring retains an aggregate advantage over
matched structural scaling and a simple alternative source-selection rule.

The original unstandardized DCT control reaches only 37.5\% coverage and
10.6\% certification.  Its difference from the standardized control shows
that the original comparison combined source scoring with a change of metric;
it cannot isolate the value of source outcomes.  Raw Sobol reaches 80.2\%
coverage and 30.0\% certification with penalized loss 0.337.


\begin{figure}[t]
  \centering
  \includegraphics[width=0.90\textwidth]{review_primary.pdf}
  \caption{Aggregate feasible coverage and independently certified deployment.
  Bars aggregate task--resolution cells; inferential comparisons are paired
  within resolution and do not treat the three grids as independent tasks.}
  \label{fig:primary-results}
\end{figure}

Paired source-minus-standardized-DCT certification differences are 0.163
(stratified bootstrap 95\% CI $[.119,.206]$) at $d=200$ and 0.156
($[.113,.200]$) at both $d=1000$ and $d=10000$.  Against the source-greedy
portfolio they are 0.113 ($[.044,.181]$) at $d=200$ and $d=1000$, and 0.081
($[.012,.150]$) at $d=10000$.  Each interval resamples the 20 paired tasks
within each of the eight fixed regimes; the resolutions are analyzed
separately.  The Online Supplement reports paired wins, losses, and ties.
These supplementary estimates describe corrected comparisons on reused tasks.
Grid refinement preserves the aggregate advantage under this task law.

\subsection{First-center attribution and prospective confirmation}

The diagnostic factorial comparison isolates the source-selected first center.
At $d=1000$, source-best with structural distance alone certifies 75/160
true-feasible deployments, compared with 74/160 for the full atlas and 49/160
for either medoid design.  Appending ranks to a fixed medoid does not recover
the gain.  Structural-only and augmented source-best designs have the same
feasible coverage at every resolution; the certification difference is one
task per resolution.  The Online Supplement gives the complete factorial and
raw-profile-distance controls.

Table \ref{tab:confirmation} reports the separate prospective experiment on
240 new latent tasks, three new source families and three new public libraries.
The atlas certifies 56/240 true-feasible deployments, compared with 24/240 for
standardized generic maximin: a paired difference of 13.3 percentage points
(95\% stratified bootstrap interval $[9.6,17.1]$).  The three family-specific
differences are 26.3, 10.0 and 3.8 points.  Source-best with structural distance
alone certifies 54/240, so the rank-distance increment is only 0.8 points
($[0,2.1]$).  The atlas exceeds the matched raw-$L^2$ design by 4.2 points
($[1.7,6.7]$), supporting a limited role for the structural metric.

The same-source portfolio certifies 49/240; the atlas-minus-portfolio
difference is 2.9 points ($[-2.1,7.9]$).  Its feasible coverage is 212/240,
compared with 155/240 for the atlas, and its mean penalized loss is 0.156
versus 0.414.  Each of the six designs records one false certificate.  Thus
the new experiment confirms a source-initialization advantage over generic
maximin, but does not confirm atlas superiority over a portfolio.  The
postdecision first-center audit finds 106/240 truly feasible source-best
centers and 83/240 medoids; it is a diagnostic, not an input to selection.
The predefined 32- and 128-profile sensitivities preserve positive paired
certification differences on the first new family (15.0 and 13.8 points),
without establishing a monotone library-size effect.

\subsection{Does structural coverage improve the remaining nine policies?}

A postdecision diagnostic fixes the same source-selected first policy and
replaces the other nine by a uniform sample without replacement from the
remaining 63 library members.  It retains all 240 tasks, original target
observation streams, ten-point budget, shortlist and verifier.  The random
remainder certifies 53/240 true-feasible deployments, versus 54/240 for
source-first structural coverage and 56/240 for the full atlas.
The paired certification differences are 0.4 percentage points
(95\% stratified bootstrap interval $[-3.3,4.2]$) and 1.3 points
($[-2.5,5.0]$), respectively.  This diagnostic does not establish a
certification advantage of structural coverage over random remainder
selection.  Initial feasible coverage is 167/240 for the random remainder
and 155/240 for either geometric design.

A fixed-first-center portfolio control changes ten ordered designs and
reuses 230 identical ordered designs.  It retains 212/240 feasible coverage
and 49/240 true certificates.  A decision-chain reconstruction explains why
these endpoints disagree: the portfolio's three-member shortlist contains a
truly feasible policy in 177 tasks, losing 35 design hits, whereas the atlas
retains all 155 hits.  Conditional on the original noisy shortlists, analytic
expectations over fresh verification streams are 54.07 true certificates for
the fixed-first portfolio and 49.72 for the atlas.  These expectations explain
verification variability; they do not replace observed outcomes.
The evidence supports source initialization over generic selection, while
leaving the incremental value of the covering rule unresolved.  The Online
Supplement reports this later diagnostic separately from the prospective
confirmation.

\subsection{Where source scoring fails}

The aggregate advantage is not uniform.  Table \ref{tab:regime-results} and
Figure \ref{fig:regime-results} preserve every adverse regime.  Raw Sobol
outperforms source scoring in aligned low frequency and irregular grids;
natural structure is slightly better under sparse high-frequency activity;
and neither method certifies useful policies under full misspecification.
The matched comparisons localize that advantage.  At $d=1000$, 24 of the
25 net certificates gained over standardized DCT come from coordinate
permutation (18) and growing rank (6).  Piecewise smoothness ties at 19/20,
and frequency shift slightly favors standardized DCT, 13/20 versus 12/20.
Against the portfolio, the largest gains occur in frequency shift and growing
rank.  In the sparse-high-frequency regime
its feasible coverage falls from 20/20 tasks at $d=200$ to 0/20 at $d=10000$,
despite unchanged target budget.  The portfolio covers 20/20 and certifies
6/20 of these $d=10000$ tasks, whereas both source-scored and standardized
maximin cover and certify none.  Geometric diversity is useful on the mixture
without being uniformly preferable to source-based greedy selection.


\begin{figure}[t]
  \centering
  \includegraphics[width=0.92\textwidth]{review_regime.pdf}
  \caption{Source-scored atlas minus each supplementary control in certified
  success, by regime and grid resolution.  Entries are percentage-point
  differences over 20 paired tasks.}
  \label{fig:regime-results}
\end{figure}

The source-scored and standardized DCT arms record no false certificates in
480 task--resolution cells each.  Unstandardized DCT, the source-greedy
portfolio, and raw Sobol record three, three, and two, respectively.
Target-only functional SCBO at
$N=13$ certifies 15 of 480 primary cells with three false certificates, all in
the frequency-support-shift regime.  Increasing its cosine rank does not remove
the general coverage deficit.

The aligned-low-frequency reversal has a concrete geometric explanation.
The supplementary geometry diagnostic reconstructs each fixed design and accesses
the hidden target safe center only afterward.  Across the 600 design points
per arm, raw Sobol has median latent-center distance 0.0343 and median non-DC
low-frequency energy 0.0258, compared with 0.0813 and 0.2292 for source
scoring.  Its pointwise true-feasible rate is 94.7\%, versus 10.0\% for both
source and the original unstandardized DCT.  High-dimensional independent coordinates average
toward a near-constant profile with DC coefficient .5; the aligned generator's
safe center happens to occupy that region.  This after-the-fact diagnostic
explains the reversal, but it uses hidden geometry and cannot be used to
choose a method or change an outcome.



\subsection{Schema and sensitivity}

Declared profile ordering is valuable information.  With domain-blind source
weights, removing the declared schema reduces certified success from 46.2\%
to 35.0\% (Table \ref{tab:schema-results}).  Descriptor-conditioned weighting
reduces both declared- and blind-schema performance and is not used in the
final method.
Source scoring therefore does not need the target family label, but it does
rely on an interpretable ordered profile map.


Sensitivity is nonmonotone, which argues against treating source volume as a
free advantage.  Raising profiles per source task from 64 to 128 reduces
certified success from 46.2\% to 38.8\%; using one, three, or five replications
gives 45.0\%, 46.2\%, and 45.6\%.  Four source tasks give 45.0\%, not more than
two.  Retaining 16 frequencies reaches 48.1\%, while four reaches 43.1\%.
Increasing $n_0$ from 10 to 20 yields only 46.9\% certification, while $n_0=5$
yields 42.5\%.  The first-center safety weight is more consequential: .25
reduces certification to 32.5\%, whereas .5 and .75 give 46.2\% and 45.0\%.
The Online Supplement displays the source-budget axes.  The
frequency-penalty sweep is algebraically redundant because column
standardization cancels positive frequency scaling; it does not demonstrate
robustness to an effective tuning parameter.



\subsection{Cost changes the ranking}

At the source method's primary budget, mean unamortized all-in calls are about
559 across the randomized law, while amortization over 20 targets reduces the
source contribution from 384 to 19.2 calls per target.  Standardized DCT uses
201.2 all-in calls and the portfolio uses 592.0 before amortization or 227.2
at 20 targets.  Verification dominates their ten target-search calls.
Figure \ref{fig:cost-results} makes both accounting
views explicit.

\begin{figure}[t]
  \centering
  \includegraphics[width=0.94\textwidth]{review_cost.pdf}
  \caption{Mean source, target-search, and verification calls under
  unamortized and 20-target accounting.}
  \label{fig:cost-results}
\end{figure}

When the 384 source calls are instead given to target-only search, the source
method is not best (Table \ref{tab:equal-cost}).  It certifies 46.2\% of 160
tasks.  Unstandardized DCT and raw Sobol each certify 53.8\%, natural structure 53.1\%,
and target-only functional SCBO 58.1\%.  The latter includes one failed cell as
an unsuccessful outcome.  Mean realized all-in costs are already close because
methods consume different verification effort.  The source archive is thus an
amortized multi-target investment, not a one-target total-cost win.


Using primary mean call counts, the archive breaks even against standardized
DCT at 15 targets when aggregating resolutions, or 15, 16, and 16 targets
at the individual resolutions.  The earlier thresholds of 7 against
unstandardized DCT, 9 against natural structure, and 11 against raw Sobol
depend on those comparators.  Both source-based rules pay 384 archive calls;
the atlas consumes less verification effort than the portfolio on average.
These are descriptive cost thresholds.

Table \ref{tab:outcome-adjusted-cost} converts amortized calls into calls per
true-feasible certificate.  With the archive shared by 20 targets, source
scoring uses 411.3 calls per certificate, versus 639.5 for standardized DCT,
612.7 for the portfolio, 1143.2 for natural structure, and 709.4 for raw
Sobol.  This descriptive ratio
does not price unsafe deployment, abstention, or objective quality.


\subsection{External Energy result}

Energy does not establish a robust certification or archive cost advantage
under the corrected physical model.  Table \ref{tab:energy-results} reports
57/90 certified source runs and 54/90 for unstandardized DCT maximin.  Standardized DCT, the
same-source portfolio, random low frequency, natural constants, raw Sobol
and functional SCBO certify 49, 59, 56, 52, 43 and 51 runs.  All eight designs
record zero false certificates.  The source-minus-unstandardized-DCT
region-balanced difference is $0.027$ (descriptive bootstrap 95\%
interval $[-.040,.120]$).  Against standardized DCT it is 0.060
($[.000,.140]$), and against the portfolio $-0.013$ ($[-.040,.000]$).
The small aggregate certification gain is unresolved across geographic
regions and does not establish superiority over the portfolio.

Of 421 certificates in the corrected eight-arm matrix, 374 satisfy the
chronological-block criterion, all 421 satisfy the nonoverlapping-window
criterion, and 374 satisfy both.  These are descriptive calendar-stability
results.  The Online Supplement reports each design separately.

Actual archive reuse across the 18 markets does not rescue the external
cost comparison.  Averaging the five repeats within each market, the source
arm consumes 6842 cumulative calls, including 3840 source calls charged once
across five regional archives, and yields 11.4 expected certified deployments.
Unstandardized DCT consumes 3258 calls and yields 10.8.  Calls per certified
success are therefore 600.2 and 301.7.  These observed reuse counts replace a
hypothetical twenty-deployment allocation for this comparison.

\subsection{Public-battery development diagnostic}

A separate retrospective replay on three 2024 seasonal tasks tests whether
source initialization survives simple policy controls.  Source-first
structural search and source-free constant-only structural search both
certify 15/15 runs across five algorithm repeats per season, with zero
finite-population false certificates and the same target query count.
Averaging repeats within each season gives 279 target queries over the three
tasks.  The source design additionally pays 44,968 source-window queries;
the constant-only design pays none.  Testing only the source-selected first
policy also certifies 15/15, using 243 target queries but retaining the same
source cost.

Paid enumeration of all 154 library policies over all 73 windows per task
requires 33,726 queries, below the source design's 45,247 source-inclusive
queries.  Restricting enumeration to the 100 constant pairs requires 21,900
queries and selects the same complete-cost optimum in each season.  This
diagnostic closes the development query-efficiency claim for source
initialization in this case.  It reports a finite-library adaptation rather
than a run of the full rank-augmented atlas.  The Online Supplement gives the
sampling law, accounting and policy selections.

\subsection{Verifier power and native transfer}

The verifier is valid but intentionally conservative.  At 80 all-success
replications, certification power is only 1.7\% at the threshold $p=.95$,
13.2\% at $.975$, and 44.8\% at $.99$ (Figure
\ref{fig:verifier-power}).  Therefore the gap between 91.7\% feasible coverage
and 47.3\% certified deployment is expected; an empty certificate set is not
evidence that search found no feasible point.  A Clopper--Pearson threshold
rule is reported as sensitivity, not substituted after observing outcomes.

\begin{figure}[t]
  \centering
  \includegraphics[width=0.86\textwidth]{review_verifier_power.pdf}
  \caption{Exact fixed-budget certification power.  The preregistered
  all-success rule is familywise valid but weak near the chance threshold.}
  \label{fig:verifier-power}
\end{figure}

In the native transfer matrix, RGPE certifies 30/60 tasks and MetaBO 28/60;
the other six methods range from 0/60 to 14/60 (Online Supplement,
subsection ``Native end-to-end transfer'').  All eight certify 0/20 Queue tasks.  These
results show that ordinary end-to-end transfer remains difficult, but they do
not establish a broad SOTA ranking from three historical synthetic families.
Their v69 verification protocol differs from the primary exact-binomial
protocol, so certification totals should not be compared across these matrices.

Finally, cumulative factor-HVD substantially improves held-out log-variance
RMSE and variance-shape correlation relative to pooled variance, yet does not
improve feasible recovery, regret, false certification, or verification cost
under matched designs.  It is therefore retained only as a calibration
diagnostic in the Online Supplement, Section EC.3, not as a main optimization
contribution.
